% =============================================================================
%  GARAH — Dossier de présentation
%  Compilable tel quel sur Overleaf (moteur : pdfLaTeX)
% =============================================================================
\documentclass[11pt,a4paper]{article}

\usepackage[T1]{fontenc}
\usepackage[utf8]{inputenc}
\usepackage[french]{babel}
\usepackage{lmodern}
\usepackage[a4paper,margin=2.4cm]{geometry}
\usepackage{graphicx}
\usepackage{xcolor}
\usepackage{booktabs}
\usepackage{array}
\usepackage{tabularx}
\usepackage{enumitem}
\usepackage{titlesec}
\usepackage{tcolorbox}
\usepackage{fancyhdr}
\usepackage{hyperref}

% --- Identité visuelle -------------------------------------------------------
\definecolor{marque}{HTML}{31AEF3}   % le bleu GARAH
\definecolor{encre}{HTML}{0F172A}
\definecolor{gris}{HTML}{64748B}
\definecolor{fondclair}{HTML}{F1F5F9}

\hypersetup{linkcolor=marque, urlcolor=marque, citecolor=marque}

\titleformat{\section}{\Large\bfseries\color{marque}}{\thesection.}{0.6em}{}
\titleformat{\subsection}{\large\bfseries\color{encre}}{\thesubsection.}{0.6em}{}
\titlespacing*{\section}{0pt}{1.6em}{0.7em}

\pagestyle{fancy}
\fancyhf{}
\lhead{\small\color{gris}GARAH — Dossier de présentation}
\rfoot{\small\color{gris}\thepage}
\renewcommand{\headrulewidth}{0.4pt}

\newtcolorbox{encadre}[1]{
  colback=fondclair, colframe=marque, boxrule=0pt, leftrule=3pt,
  arc=2pt, left=10pt, right=10pt, top=8pt, bottom=8pt,
  title={\bfseries\color{marque}#1}, coltitle=marque,
  fonttitle=\bfseries, colbacktitle=fondclair, titlerule=0pt
}

\newcommand{\chiffre}[2]{%
  \begin{minipage}[t]{0.30\textwidth}
    \centering
    {\Huge\bfseries\color{marque}#1}\\[0.3em]
    {\small\color{gris}#2}
  \end{minipage}%
}

\setlist[itemize]{leftmargin=1.2em, itemsep=0.25em, topsep=0.4em}

% =============================================================================
\begin{document}

% --- Couverture --------------------------------------------------------------
\begin{titlepage}
\centering
\vspace*{3cm}

{\Huge\bfseries\color{marque} GARAH}\\[0.6em]
{\Large\color{encre} Vendre, servir et acheminer}\\[0.3em]
{\Large\color{encre} entre le Cameroun et la Centrafrique}\\[3cm]

\begin{minipage}{0.8\textwidth}
\centering\large\color{gris}
Une plateforme unique qui relie le commerce en ligne,
la relation client et le transport de marchandises
sur le corridor Douala~--~Bertoua~--~Bangui.
\end{minipage}

\vfill
{\color{gris}Dossier de présentation aux partenaires \\[0.3em]
\today}
\end{titlepage}

\tableofcontents
\newpage

% =============================================================================
\section{En une page}
% =============================================================================

\subsection*{Le problème}

Acheter un produit à Douala et le recevoir à Bangui relève aujourd'hui du
système D. Le commerçant prend la commande par téléphone ou par messagerie,
note le paiement sur un cahier, confie le colis à un transporteur, et
répond au client « il est parti » sans pouvoir dire où il se trouve.

Quand un colis se perd, personne ne sait à quelle étape. Quand un client
conteste, personne ne peut rouvrir l'historique. Quand un marchand partenaire
demande son argent, le calcul se refait à la main.

\subsection*{La réponse}

GARAH réunit dans un seul système ce qui est aujourd'hui éparpillé entre un
cahier, une messagerie et un carnet de transporteur :

\begin{itemize}
  \item une \textbf{boutique en ligne} où le client choisit, négocie et paie
        par Mobile Money~;
  \item un \textbf{service client} intégré, où la discussion peut se conclure
        par une vente~;
  \item un \textbf{suivi de colis} étape par étape, consultable par le client
        avec un simple numéro reçu par SMS~;
  \item une \textbf{comptabilité automatique} des sommes dues à chaque
        marchand partenaire.
\end{itemize}

\subsection*{L'état d'avancement}

\vspace{0.6em}
\noindent
\chiffre{9 / 11}{lots fonctionnels\\livrés}
\hfill
\chiffre{100\,\%}{du back-office\\opérationnel}
\hfill
\chiffre{En ligne}{plateforme déployée\\et accessible}
\vspace{1.2em}

Le cœur du système fonctionne et tourne en production. Ce qui reste à
construire est la boutique vue par le client final~: le moteur est en place,
il faut lui poser la carrosserie.

\newpage
% =============================================================================
\section{Le marché et le besoin}
% =============================================================================

\subsection{Un corridor commercial réel, mal outillé}

L'axe \textbf{Douala~--~Bertoua~--~Bangui} est une route commerciale active.
Des marchandises y circulent tous les jours~: produits manufacturés, textiles,
équipements, pièces détachées. Mais la chaîne repose sur la confiance
interpersonnelle et sur des traces écrites dispersées.

\begin{encadre}{Ce que coûte l'absence d'outil}
\begin{itemize}
  \item \textbf{Le litige sans preuve.} Sans historique daté, un désaccord se
        règle à la parole contre la parole. Le commerçant perd le client, ou
        rembourse à tort.
  \item \textbf{L'argent qui s'évapore.} Une commission oubliée sur une vente
        ne se remarque qu'en fin de mois, quand elle n'est plus récupérable.
  \item \textbf{Le client qui rappelle.} « Où est mon colis ? » occupe un
        employé plusieurs fois par jour, pour une réponse que le système
        devrait donner tout seul.
\end{itemize}
\end{encadre}

\subsection{Trois publics, trois besoins distincts}

\begin{center}
\begin{tabularx}{\textwidth}{@{}l>{\raggedright\arraybackslash}X@{}}
\toprule
\textbf{Public} & \textbf{Ce qu'il attend} \\
\midrule
Le \textbf{client final} &
Voir ce qui est disponible, connaître le prix réel selon la quantité, payer
avec son téléphone, et savoir où est sa commande sans appeler personne. \\
\addlinespace
Le \textbf{marchand partenaire} &
Confier ses produits à une vitrine qui a du trafic, et être payé sur un
décompte qu'il peut vérifier ligne par ligne. \\
\addlinespace
L'\textbf{entreprise} &
Piloter les stocks, le transport, les réclamations et l'argent depuis un
seul écran, avec la trace de qui a fait quoi. \\
\bottomrule
\end{tabularx}
\end{center}

\newpage
% =============================================================================
\section{Ce que la plateforme sait faire}
% =============================================================================

\subsection{Vendre}

\begin{itemize}
  \item \textbf{Un catalogue à déclinaisons.} Un même article existe en
        tailles, couleurs et conditionnements~; chaque version a son propre
        prix et son propre stock.
  \item \textbf{Des prix dégressifs.} Le tarif dépend de la quantité commandée~:
        un prix de 1 à 6 pièces, un autre au-delà. C'est le fonctionnement
        réel du commerce de gros, pas une option.
  \item \textbf{La négociation.} Le client discute avec un conseiller, et une
        proposition de prix acceptée devient directement une commande. Le
        marchandage n'est pas contourné~: il est outillé.
  \item \textbf{Le paiement Mobile Money.} MTN MoMo et Orange Money, les
        moyens réellement utilisés au Cameroun.
\end{itemize}

\begin{encadre}{Un choix qui protège l'entreprise}
Les montants d'une commande sont \textbf{figés au moment de l'achat} —
désignation, prix, taxe, commission, frais de livraison. Changer un tarif
aujourd'hui ne réécrit pas une facture de mars. Sans cette règle, toute la
comptabilité passée deviendrait fausse à chaque mise à jour de prix.
\end{encadre}

\subsection{Servir}

\begin{itemize}
  \item \textbf{Une file d'attente partagée.} Les conversations en attente sont
        visibles de toute l'équipe~; le premier conseiller disponible prend le
        dossier.
  \item \textbf{Des réclamations et des retours traités séparément.} Une
        réclamation est une plainte qu'un humain tranche~; un retour est une
        marchandise qui revient physiquement. Ce ne sont ni les mêmes gestes,
        ni les mêmes personnes.
  \item \textbf{Un retour en cinq étapes} — accepté, reçu, contrôlé, validé,
        clôturé. Un retour accepté n'est pas un retour reçu~: confondre les
        deux, c'est rembourser une marchandise qui n'est jamais rentrée.
  \item \textbf{Une évaluation après clôture seulement.} Demander un avis sur
        un problème non résolu ne mesure pas la qualité du service, il mesure
        l'agacement.
\end{itemize}

\subsection{Acheminer}

\begin{itemize}
  \item \textbf{Des points de transit et de récupération} le long du corridor.
        Le client choisit où il vient chercher sa commande.
  \item \textbf{Des itinéraires réutilisables}, qui annoncent un délai sans
        jamais l'imposer~: une route coupée, cela arrive, et le système doit
        pouvoir enregistrer la réalité.
  \item \textbf{Un suivi par événements datés} — « réceptionné à Bertoua le
        12/03 à 14~h par David » — et non par un simple statut qui s'écrase.
        Le parcours reste lisible même après la livraison.
  \item \textbf{Un suivi public.} Le client entre le numéro reçu par SMS et
        voit où en est son colis, sans compte et sans mot de passe.
  \item \textbf{Un comptoir de remise.} L'agent saisit le code, \emph{voit}
        les colis qu'il désigne, puis confirme. Voir avant de remettre évite
        de rendre la mauvaise marchandise.
\end{itemize}

\subsection{Compter}

\begin{itemize}
  \item \textbf{Un grand livre par marchand.} Chaque vente écrit une ligne,
        commission comprise. Le solde est la somme de ces lignes, jamais un
        total conservé à part — un total conservé finit toujours par mentir.
  \item \textbf{Des règlements en deux temps} — préparer, puis confirmer. La
        dette ne baisse qu'au moment où l'argent part réellement, et la
        référence du virement est obligatoire~: c'est la seule preuve six mois
        plus tard.
  \item \textbf{Des corrections par ajout.} Une écriture fausse s'annule par
        une autre écriture, avec son motif et son auteur. Un grand livre ne se
        réécrit pas.
  \item \textbf{Des statistiques de vente et d'audience} — chiffre d'affaires,
        articles vendus, fiches consultées, produits qui montent — exportables
        en PDF, Word et Excel.
\end{itemize}

\newpage
% =============================================================================
\section{Ce qui distingue GARAH}
% =============================================================================

\subsection{Un système pensé pour le litige, pas seulement pour la vente}

La plupart des outils de commerce enregistrent ce qui se passe quand tout va
bien. GARAH est construit autour de ce qui se passe quand ça va mal —~parce
que c'est là que l'argent se perd.

\begin{center}
\begin{tabularx}{\textwidth}{@{}>{\raggedright\arraybackslash}X>{\raggedright\arraybackslash}X@{}}
\toprule
\textbf{Le réflexe habituel} & \textbf{Le choix de GARAH} \\
\midrule
Une colonne « statut » qu'on met à jour. &
Un \textbf{journal d'événements} : on garde chaque étape, pas seulement la
dernière. \\
\addlinespace
Supprimer ce qui ne sert plus. &
\textbf{Désactiver}, jamais supprimer : des factures et des livraisons
pointent vers ces lignes. \\
\addlinespace
Un total de solde stocké quelque part. &
Un solde \textbf{recalculé} à partir des écritures, donc toujours juste. \\
\addlinespace
Un bouton qui échoue et affiche « erreur ». &
L'écran dit \textbf{ce qui manque avant le clic}. \\
\bottomrule
\end{tabularx}
\end{center}

\subsection{Une traçabilité qui tient devant un tiers}

\begin{itemize}
  \item \textbf{Qui a fait quoi} : un journal d'audit enregistre chaque
        modification interne.
  \item \textbf{Ce que le client a vécu} : son parcours est mesuré séparément.
  \item \textbf{Ce qui touche à la sécurité} : connexions et alertes sont
        journalisées à part.
\end{itemize}

Ces trois journaux ne se mélangent pas, précisément pour qu'une question
—~« qui a annulé cette commande ? »~— trouve une réponse nette.

\subsection{Des droits fins, et une séparation des pouvoirs}

Le système distingue quatre niveaux d'accès, et \textbf{188 autorisations}
distinctes que l'on assemble en profils métier (« gestionnaire de catalogue »,
« agent de comptoir », « service client »).

\begin{center}
\begin{tabularx}{\textwidth}{@{}l>{\raggedright\arraybackslash}X@{}}
\toprule
\textbf{Niveau} & \textbf{Rôle} \\
\midrule
Super administrateur & Décide de ce que le système sait faire, et qui surveille qui. \\
Administrateur & Décide qui travaille, et sur quoi. \\
Chef de service & Encadre son équipe. \\
Responsable & Exécute, selon les profils qu'on lui a donnés. \\
\bottomrule
\end{tabularx}
\end{center}

\begin{encadre}{Pourquoi cela compte pour un investisseur}
Un agent de service client ne peut pas déclencher un remboursement. Trancher
un litige en faveur du client et rendre l'argent sont deux décisions
séparées, avec deux autorisations différentes. C'est ce qui empêche qu'une
seule personne puisse, seule, faire sortir de l'argent.
\end{encadre}

\subsection{Conçu pour les conditions locales}

\begin{itemize}
  \item \textbf{Léger sur les connexions lentes.} L'application se charge en
        une centaine de kilo-octets~; les fonctions lourdes ne sont
        téléchargées qu'au moment où on les utilise.
  \item \textbf{Trois langues prévues} dès l'origine — français, anglais,
        sango.
  \item \textbf{Le franc CFA} comme monnaie de référence.
  \item \textbf{Le Mobile Money} comme moyen de paiement principal, et non
        comme option secondaire ajoutée après coup.
\end{itemize}

\newpage
% =============================================================================
\section{Où en est le projet}
% =============================================================================

\subsection{Ce qui est livré et fonctionne}

\begin{center}
\begin{tabularx}{\textwidth}{@{}>{\raggedright\arraybackslash}Xc@{}}
\toprule
\textbf{Domaine} & \textbf{État} \\
\midrule
Catalogue, déclinaisons et grilles de prix & Livré \\
Commandes et paiements & Livré \\
Stocks et mouvements & Livré \\
Logistique : lieux, itinéraires, expéditions, suivi, comptoir & Livré \\
Réclamations et retours & Livré \\
Service client et négociation & Livré \\
Finance marchand : grand livre, règlements, commissions & Livré \\
Clients, surveillance et audit & Livré \\
Statistiques et exports & Livré \\
\addlinespace
Boutique vue par le client final & En cours \\
\bottomrule
\end{tabularx}
\end{center}

\subsection{Ce que ces chiffres veulent dire}

\vspace{0.8em}
\noindent
\chiffre{64}{tables de données\\structurées}
\hfill
\chiffre{188}{autorisations\\distinctes}
\hfill
\chiffre{345}{contrôles automatiques\\au vert}
\vspace{1.2em}

Les \textbf{345 contrôles automatiques} méritent une explication. Ce sont des
vérifications que la machine rejoue à chaque modification du système. Elles
répondent à une question simple~: « est-ce que ce qui marchait hier marche
encore aujourd'hui ? »

\begin{encadre}{Pourquoi c'est un argument financier}
Sur un logiciel sans ces contrôles, chaque évolution risque de casser
silencieusement autre chose, et le défaut n'apparaît que chez le client. Le
coût de maintenance grimpe avec le temps, jusqu'à ce que plus personne n'ose
toucher au système.

Ici, une régression est détectée \textbf{en quelques minutes}, avant d'être
mise en ligne. C'est ce qui rend le projet évolutif sur la durée, et non
seulement fonctionnel aujourd'hui.
\end{encadre}

\subsection{Une décision, une trace}

Le projet tient un \textbf{journal des décisions} — 31 entrées à ce jour.
Chacune écrit ce qui a été choisi, pourquoi, et ce qu'on y perd.

C'est inhabituel, et c'est délibéré~: dans deux ans, quand quelqu'un voudra
modifier un comportement qui semble arbitraire, il trouvera la raison au lieu
de rouvrir un problème déjà résolu. Le savoir ne repose pas sur la mémoire
d'une personne.

\newpage
% =============================================================================
\section{Les fondations techniques}
% =============================================================================

\emph{Cette section est volontairement brève. Elle sert à établir que les
choix sont solides et courants, pas exotiques.}

\subsection{Ce sur quoi c'est construit}

\begin{center}
\begin{tabularx}{\textwidth}{@{}l>{\raggedright\arraybackslash}X@{}}
\toprule
\textbf{Élément} & \textbf{Choix, et ce que cela garantit} \\
\midrule
Serveur applicatif & \textbf{Java / Spring Boot} — la technologie des banques
et des assurances. Compétences largement disponibles. \\
\addlinespace
Base de données & \textbf{PostgreSQL} — libre, éprouvée, sans coût de licence.
Hébergée en Europe. \\
\addlinespace
Interfaces & \textbf{Angular} — maintenu par Google, adapté aux applications
de gestion lourdes. \\
\addlinespace
Hébergement & \textbf{Microsoft Azure}, région France. Serveur toujours actif,
sans temps de réveil. \\
\addlinespace
Fichiers et photos & Stockage objet dédié, avec adresses signées à durée
limitée. \\
\bottomrule
\end{tabularx}
\end{center}

\subsection{Deux garanties d'exploitation}

\begin{itemize}
  \item \textbf{Aucune dépendance à un fournisseur unique.} Toutes les briques
        sont standard et remplaçables. Le système a déjà été déplacé d'un
        hébergeur à un autre en cours de route, sans réécriture.
  \item \textbf{Mise en ligne automatisée.} Chaque modification est vérifiée
        par les 345 contrôles, puis déployée sans intervention manuelle. Une
        version défectueuse n'atteint pas la production.
\end{itemize}

\newpage
% =============================================================================
\section{La suite}
% =============================================================================

\subsection{Prochaine étape : la boutique client}

C'est le dernier grand chantier. Le serveur expose déjà tout ce dont elle a
besoin~: catalogue, panier, commande, paiement, suivi, réclamations,
conversations. Il s'agit de construire les écrans que verra le client final.

\begin{center}
\begin{tabularx}{\textwidth}{@{}l>{\raggedright\arraybackslash}X@{}}
\toprule
\textbf{Écran} & \textbf{Fonction} \\
\midrule
Vitrine et catégories & Le catalogue public. \\
Fiche produit & Choix de la déclinaison, prix selon la quantité. \\
Panier et commande & Adresse, point de récupération, récapitulatif. \\
Paiement & MTN MoMo et Orange Money. \\
Mes commandes et suivi & Où en est chaque colis. \\
Réclamations et conversation & Le service client, côté client. \\
\bottomrule
\end{tabularx}
\end{center}

\subsection{Ce qu'un partenariat permettrait d'accélérer}

\begin{itemize}
  \item \textbf{Finir la boutique client} et ouvrir la plateforme au public.
  \item \textbf{Recruter les marchands partenaires} du corridor, et les
        accompagner à la mise en ligne de leur catalogue.
  \item \textbf{Équiper les points de récupération} — le comptoir de remise
        existe, il faut le déployer sur le terrain.
  \item \textbf{Conventionner les opérateurs Mobile Money} pour des conditions
        adaptées au volume.
\end{itemize}

\vspace{1.5em}
\begin{encadre}{Ce qui est déjà acquis}
Le risque technique principal d'un projet de cette nature —~construire un
système fiable qui tienne la comptabilité, la logistique et le service
client ensemble~— est \textbf{derrière nous}. Le système existe, il tourne,
il est vérifié en continu et documenté.

Ce qui reste relève de l'exécution commerciale et du déploiement terrain.
\end{encadre}

\vfill
\begin{center}
\color{gris}\small
GARAH — Document de présentation aux partenaires \\
Ce document décrit l'état du projet à la date indiquée en couverture.
\end{center}

\end{document}
